\chapter{The Idea of Frequency}

\section{What we are really after}

We have spent the previous chapters building two ideas with our bare hands. First,
a signal is a function: a number that changes as time goes by, or as you move across
a sensor. Second, a sinusoid is the most special function there is, because it is
the shape that survives being differentiated, integrated, delayed, and pushed through
any linear time-invariant system without changing its \emph{shape} -- only its size
and its timing. In this chapter we put those two ideas together and ask the question
that organizes the whole rest of the book:

\begin{focusbox}
What does it mean to say a signal ``contains'' a frequency? And why is it so
powerful to describe a signal not by its values moment to moment, but by the
mixture of frequencies it is built from?
\end{focusbox}

This is not a small question. The answer -- that \emph{any} reasonable signal
can be written as a sum of pure sinusoids, and that this sum is unique -- is the
Fourier idea, and it took a century of argument to make precise. We are not going to
prove it yet. We are going to \emph{see} it. By the end of the chapter you will be
able to look at a wiggly waveform and have a real intuition for what its spectrum
must look like, and look at a spectrum and imagine the waveform. That two-way
fluency is the single most useful skill in all of signal processing.

\begin{goalbox}
The wrist controller we are building toward lives or dies on this idea. A finger
tap, a wrist flick, a tremor, the 60~Hz hum the wall socket induces in the EMG
leads, the steady gravity vector the accelerometer reports -- each of these shows
up as energy at a particular \emph{frequency} or band of frequencies. ``Is the user
tapping rhythmically at 4~Hz, or is that just arm sway at 0.5~Hz?'' is a question
about the spectrum, not about any single sample. Learning to think in frequency is
learning the native language of the device.
\end{goalbox}

\section{Frequency as a rate of repetition}

Start with the plainest possible meaning. A signal is \emph{periodic} with period $T$
if it repeats exactly after $T$ seconds:
\[
  x(t+T) = x(t) \quad \text{for all } t.
\]
The smallest such $T$ is \emph{the} period. The \emph{frequency} is just how many
times that repetition happens per second:
\[
  f = \frac{1}{T} \qquad \text{[cycles per second, or hertz, Hz]}.
\]
A sinusoid that completes one full cycle every $0.25$~s repeats $4$ times a second,
so $f = 4$~Hz. Double the frequency and the pattern is squeezed to half the width
in time; halve the frequency and it stretches to twice the width. That inverse
relationship between time and frequency -- squeeze one, stretch the other -- is a
theme we will meet again and again, all the way to the uncertainty principle in
Chapter~10.

\begin{intuitionbox}
Frequency is a \emph{rate}. It answers ``how fast does this pattern recur?'' Low
frequency means slow, lazy, broad wiggles. High frequency means fast, tight,
needle-like wiggles. Everything else in this chapter is an elaboration of that one
sentence, made precise enough to compute with.
\end{intuitionbox}

\subsection{Ordinary frequency versus angular frequency}

There are two units for the same physical rate, and confusing them is the single
most common bookkeeping error in the field. A pure cosine can be written either way:
\[
  x(t) = \cos(2\pi f t) = \cos(\omega t), \qquad \omega = 2\pi f.
\]
\begin{itemize}
  \item \textbf{Ordinary frequency} $f$ is measured in hertz: cycles per second.
        It is what a tuner, an oscilloscope, or a datasheet quotes. A 440~Hz tuning
        fork has $f = 440$.
  \item \textbf{Angular frequency} $\omega$ is measured in radians per second. One
        full cycle is $2\pi$ radians of phase, so $\omega = 2\pi f$. It is what the
        math wants, because the argument of $\cos$ and of $e^{\jj\theta}$ is an
        \emph{angle}.
\end{itemize}
The reason both survive is that hardware counts cycles and mathematics counts angles,
and there are $2\pi$ radians in a cycle. Keep a mental dictionary: whenever you see a
bare $\omega$, divide by $2\pi$ in your head to get the hertz you can feel.
Table~\ref{ch05:units} collects the three names.

\begin{trapbox}
The phrase ``a frequency of $\omega$'' is ambiguous unless units are stated. Worse,
the discrete world (Chapter~6 onward) uses yet a third creature: the normalized
\emph{digital} frequency in radians per \emph{sample}, often also called $\omega$,
where $\pi$ corresponds to half the sampling rate. Always ask: per second or per
sample? Hertz or radians? When in doubt, write $2\pi f t$ in full -- it is never
wrong.
\end{trapbox}

\begin{table}[H]
\centering
\begin{tabular}{@{}llll@{}}
\toprule
Name & Symbol & Units & Relation \\
\midrule
Period            & $T$      & seconds            & $T = 1/f$ \\
Ordinary freq.    & $f$      & hertz (cycles/s)   & $f = 1/T$ \\
Angular freq.     & $\omega$ & radians/second     & $\omega = 2\pi f$ \\
\bottomrule
\end{tabular}
\caption{The three names for one rate. Memorize $\omega = 2\pi f$ and $f = 1/T$.}
\label{ch05:units}
\end{table}

\section{The atoms: pure sinusoids and their three numbers}

A general single-frequency real signal needs exactly three numbers to pin down:
\[
  x(t) = A\cos(2\pi f t + \phi).
\]
\begin{itemize}
  \item $A$ -- the \emph{amplitude}: how tall the wave is, how loud, how strong.
  \item $f$ -- the \emph{frequency}: how fast it repeats.
  \item $\phi$ -- the \emph{phase}: where in its cycle it is at $t=0$, i.e.\ a
        sideways shift in time. A phase of $-\pi/2$ turns a cosine into a sine.
\end{itemize}
Three numbers, total. That economy is the whole point. A raw recorded waveform might
be ten thousand samples; if it happens to be a single sinusoid, those ten thousand
numbers carry only \emph{three} numbers' worth of information. Frequency-domain
thinking is, at bottom, the search for that hidden economy in every signal.
Figure~\ref{ch05:twosines} shows two such atoms side by side.

\begin{figure}[H]
\centering
\begin{tikzpicture}
\begin{axis}[
  width=12cm, height=5cm,
  title={Two sinusoids over one second ($2$\,Hz vs $5$\,Hz)},
  title style={font=\small},
  xlabel={time $t$ (s)}, ylabel={$x(t)$},
  domain=0:1, samples=200,
  axis lines=left, enlargelimits=false,
  xmin=0, xmax=1, ymin=-1.6, ymax=1.6,
  legend pos=outer north east, legend cell align=left,
]
\addplot[sigblue, very thick] {cos(deg(2*pi*2*x))};
\addlegendentry{$\cos(2\pi\cdot 2\,t)$, $f=2$ Hz}
\addplot[specred, very thick, dashed] {0.6*cos(deg(2*pi*5*x))};
\addlegendentry{$0.6\cos(2\pi\cdot 5\,t)$, $f=5$ Hz}
\end{axis}
\end{tikzpicture}
\caption{Two pure sinusoids over one second. The blue wave completes $2$ cycles
(period $0.5$~s); the red wave completes $5$ cycles (period $0.2$~s) and is
shorter. Higher frequency = tighter wiggle; smaller amplitude = shorter wave.}
\label{ch05:twosines}
\end{figure}

\begin{connectionbox}
In the linear-algebra companion you met orthonormal bases: a handful of perpendicular
unit vectors so that every vector is a unique weighted sum of them, and the weights
are inner products. Keep that picture warm. The Fourier claim is exactly that the
sinusoids form such a basis for the space of signals, with the inner product
$\inner{x}{y} = \int x(t)\,\overline{y(t)}\,\dd t$. ``How much of frequency $f$ does
$x$ contain?'' will literally be an inner product. We are sneaking up on that.
\end{connectionbox}

\section{Why decompose into frequencies at all?}

Before the machinery, the motivation. Why is breaking a signal into sinusoids worth
the trouble, rather than, say, breaking it into a sum of little spikes or boxes?

\begin{whybox}
\textbf{Reason 1: sinusoids are the eigenfunctions of every LTI system.} Push a
sinusoid through any linear time-invariant system -- a filter, a cable, a room, a
mechanical linkage -- and out comes \emph{the same frequency}, only rescaled in
amplitude and shifted in phase. No new frequencies are ever created. We will prove
this in Chapter~12, but the consequence is decisive: if you describe the input as a
sum of sinusoids, the system acts on each one \emph{independently and by simple scalar multiplication}
(multiply by a complex number). A horrendous integral (convolution) becomes
elementary multiplication. Spikes and boxes do not have this property; sinusoids
alone do.
\end{whybox}

\begin{whybox}
\textbf{Reason 2: it matches how the world generates and senses signals.} A plucked
string vibrates at a fundamental and its harmonics. A motor hums at its rotation
rate and multiples. The cochlea in your ear is literally a bank of frequency-tuned
hairs -- you \emph{hear} in the frequency domain. AC mains radiate at 50 or 60~Hz.
When the physics speaks in frequencies, analysis in frequencies is not a trick; it
is reading the world in its own handwriting.
\end{whybox}

\begin{whybox}
\textbf{Reason 3: it separates things that overlap in time but not in frequency.}
A drum hit and a flute note can happen at the same instant -- hopelessly tangled in
the time waveform -- yet sit in completely different frequency bands, easily separated
once you look at the spectrum. Most of practical DSP, including pulling a
finger-tap transient out of slow arm motion, is exactly this kind of separation.
\end{whybox}

\section{Building a waveform from sinusoids: the gentle path}

Here is the experiment that, more than any equation, makes the Fourier idea click.
We will build a non-sinusoidal shape -- a square wave -- by adding sinusoids, one at
a time, and watch the target emerge.

Take odd harmonics of a fundamental $f_0$, each with amplitude $1/k$:
\[
  s_N(t) = \frac{4}{\pi}\sum_{\substack{k=1\\k\ \text{odd}}}^{N}
           \frac{1}{k}\sin(2\pi k f_0 t)
         = \frac{4}{\pi}\!\left(\sin\omega_0 t + \tfrac13\sin 3\omega_0 t
           + \tfrac15\sin 5\omega_0 t + \cdots\right).
\]
One term is a lonely sine. Two terms already has the rough plateau-and-edge look
of a square. Many terms and it is unmistakably a square wave with sharp risers, as
Figure~\ref{ch05:square} shows.

\begin{figure}[H]
\centering
\begin{tikzpicture}
\begin{axis}[
  width=12cm, height=5.5cm,
  title={Square wave approximation from odd harmonics},
  title style={font=\small},
  xlabel={time $t$ (one period)}, ylabel={amplitude},
  domain=0:1, samples=200,
  axis lines=left, enlargelimits=false,
  xmin=0, xmax=1, ymin=-1.5, ymax=1.5,
  legend pos=outer north east, legend cell align=left,
]
% 1 term
\addplot[sigblue, thick] {(4/pi)*sin(deg(2*pi*x))};
\addlegendentry{1 term}
% 3 terms
\addplot[quotientorange, thick] {(4/pi)*(sin(deg(2*pi*x)) + (1/3)*sin(deg(2*pi*3*x)))};
\addlegendentry{2 terms}
% many terms
\addplot[specred, very thick] {(4/pi)*(sin(deg(2*pi*x)) + (1/3)*sin(deg(2*pi*3*x)) + (1/5)*sin(deg(2*pi*5*x)) + (1/7)*sin(deg(2*pi*7*x)) + (1/9)*sin(deg(2*pi*9*x)) + (1/11)*sin(deg(2*pi*11*x)) + (1/13)*sin(deg(2*pi*13*x)))};
\addlegendentry{7 terms}
\end{axis}
\end{tikzpicture}
\caption{A square wave emerging from odd harmonics. One sine is a poor square; add
the third and fifth and seventh harmonics (here up to the 13th) and the flat tops
and steep edges appear. The little overshoot at the jumps never goes away -- that is
the Gibbs phenomenon, revisited in Chapter~6.}
\label{ch05:square}
\end{figure}

\begin{happeningbox}
Watch what each new term \emph{does}. The fundamental sets the overall up-down
rhythm. Each higher harmonic is a faster, smaller correction that sharpens the
corners and flattens the tops. The \emph{slow} sinusoids carry the gross shape; the
\emph{fast} ones carry the fine detail and the edges. This is the deepest single
intuition about spectra: \textbf{low frequencies are gross shape, high frequencies
are sharp detail and edges.} A blurry image is one with the high frequencies removed;
a muffled voice is one with the highs rolled off; a signal with sharp transients
must have substantial high-frequency content. You can now reason about a spectrum
without computing a single integral.
\end{happeningbox}

\begin{labbox}
Type this and drag the slider in your head from 1 term to 50. The corners get
crisper; the overshoot at the edge stubbornly stays about 9\% tall. That stubborn
ear on the jump is telling you something real about representing a discontinuity
with smooth waves.
\begin{lstlisting}[language=Python]
import numpy as np, matplotlib.pyplot as plt
t = np.linspace(0, 1, 1000, endpoint=False)
def square_approx(N):
    s = np.zeros_like(t)
    for k in range(1, N+1, 2):          # odd harmonics only
        s += (1/k) * np.sin(2*np.pi*k*t)
    return (4/np.pi) * s
for N in (1, 3, 7, 49):
    plt.plot(t, square_approx(N), label=f"{(N+1)//2} terms")
plt.legend(); plt.xlabel("t"); plt.show()
\end{lstlisting}
\end{labbox}

\section{The spectrum: a signal seen edge-on}

If a signal is a sum of sinusoids, then to specify it you need only list, for each
frequency present, its amplitude and its phase. That list \emph{is} the spectrum.

\begin{definition}[Spectrum, informal]
\label{ch05:spectrum-def}
The \emph{spectrum} of a signal is the description of which frequencies it contains
and in what amount. For each frequency $f$ we record an \emph{amplitude} (or
magnitude) $A(f)$ -- how much of that frequency is present -- and a \emph{phase}
$\phi(f)$ -- the timing of that component. Plotting $A$ against $f$ gives the
\emph{magnitude spectrum}; plotting $\phi$ against $f$ gives the \emph{phase
spectrum}.
\end{definition}

Two ways to draw the very same square wave, in Figure~\ref{ch05:twoviews}:

\begin{figure}[H]
\centering
\begin{tikzpicture}
% time domain mini
\begin{axis}[
  name=tplot, width=6.2cm, height=4.2cm,
  title={Time domain $x(t)$},
  title style={font=\small},
  xlabel={$t$}, ylabel={$x(t)$},
  domain=0:1, samples=200, axis lines=left,
  xmin=0, xmax=1, ymin=-1.4, ymax=1.4, ytick={-1,0,1},
]
\addplot[sigblue, very thick] {(4/pi)*(sin(deg(2*pi*x)) + (1/3)*sin(deg(2*pi*3*x)) + (1/5)*sin(deg(2*pi*5*x)) + (1/7)*sin(deg(2*pi*7*x)) + (1/9)*sin(deg(2*pi*9*x)))};
\end{axis}
% frequency domain mini (stems)
\begin{axis}[
  at={(tplot.east)}, anchor=west, xshift=1.4cm,
  width=6.2cm, height=4.2cm,
  title={Magnitude spectrum},
  title style={font=\small},
  xlabel={harmonic $k$ ($f = k f_0$)},
  ylabel={amplitude},
  axis lines=left, xmin=0, xmax=10, ymin=0, ymax=1.45,
  xtick={1,3,5,7,9}, ytick={0,0.5,1},
  ycomb, mark=*, mark options={specred},
]
\addplot[specred, very thick] coordinates {(1,1.273) (3,0.424) (5,0.255) (7,0.182) (9,0.141)};
\end{axis}
\end{tikzpicture}
\caption{The same signal, two views. Left: the wiggle in time. Right: its line
spectrum -- a stem at each odd harmonic whose height is $\frac{4}{\pi k}$. The
spectrum is the signal ``seen edge-on'': all the time-domain complexity collapses to
a handful of dots.}
\label{ch05:twoviews}
\end{figure}

\begin{aestheticbox}
Notice the economy. The left plot needs hundreds of points to draw; the right plot
needs five dots. Same information. The spectrum is not an approximation or a summary
-- for a periodic signal it is a \emph{complete, lossless} re-encoding. This is the
quiet miracle that makes MP3, JPEG, radio, and the FFT possible: most real signals
are sparse or simple in frequency even when they look hopeless in time.
\end{aestheticbox}

\subsection{The Fourier winding machine: finding the hidden recipe}
\label{ch05:winding}

How do you actually uncover what frequencies live inside a tangled curve? The geometric
picture popularized by 3Blue1Brown makes this visceral without formulas.

Imagine taking your time signal $x(t)$ and \emph{winding it around the origin} of the 2D complex plane.
At each instant $t$, place a point at distance $x(t)$ from the origin, rotating at a variable test
frequency $f_{\text{test}}$:
\[
  z(t) = x(t)\,e^{-\jj 2\pi f_{\text{test}} t} .
\]
Now imagine this wound path is made of wire with uniform mass density, and ask: \emph{where is its center of mass?}
The center of mass is simply the time average
\[
  \mu(f_{\text{test}}) = \frac{1}{T}\int_0^T x(t)\,e^{-\jj 2\pi f_{\text{test}} t}\,\dd t .
\]
Figure~\ref{ch05:fig-winding} reveals what happens as you sweep the winding frequency $f_{\text{test}}$:

\begin{itemize}
  \item \textbf{When $f_{\text{test}}$ does not match any rhythm in $x(t)$}: The crests and troughs of the
  wave get sprayed evenly around all angles of the circle. The lobes balance each other out in pairs, and the
  center of mass sits trapped right at the origin, $\mu \approx 0$.
  \item \textbf{When $f_{\text{test}}$ matches a frequency present in $x(t)$}: Every crest of the wave arrives
  at the exact same angle on the circle! The wire bulges heavily toward one side. The center of mass leaps outward
  away from the origin. The distance of $\mu$ from the center equals half the amplitude of that frequency component,
  and its angle indicates the phase!
\end{itemize}

As you dial the winding frequency $f_{\text{test}}$ up and down, the center of mass hovers near zero almost
everywhere, only jumping outward when $f_{\text{test}}$ hits a genuine frequency component of the signal.
\textbf{The Fourier transform is nothing other than tracking the center of mass of the wound signal as a function of winding frequency.}

\begin{figure}[H]
\centering
\begin{tikzpicture}
  % Panel 1: Non-matching winding frequency (f_test = 2 Hz, signal = 3 Hz)
  \begin{scope}[xshift=0cm]
    \draw[->,gray!50] (-1.9,0) -- (1.9,0) node[right,font=\tiny]{$\mathrm{Re}$};
    \draw[->,gray!50] (0,-1.9) -- (0,1.9) node[above,font=\tiny]{$\mathrm{Im}$};
    \draw[gray!30,dashed] (0,0) circle (1);
    \draw[sigblue,thick,domain=0:360,samples=160,smooth,variable=\deg]
      plot ({-(1 + 0.8*cos(1.5*\deg))*cos(\deg)}, {-(1 + 0.8*cos(1.5*\deg))*sin(\deg)});
    \fill[specred] (0,0) circle (2.2pt) node[below right=1pt,font=\footnotesize\bfseries]{$\mu \approx 0$};
    \node[anchor=north,font=\small] at (0,-2.1) {(a) $f_{\text{test}} = 2$\,Hz $\neq 3$\,Hz};
    \node[anchor=north,font=\scriptsize,gray!70!black] at (0,-2.5) {symmetric lobes cancel};
  \end{scope}

  % Panel 2: Matching winding frequency (f_test = 3 Hz, signal = 3 Hz)
  \begin{scope}[xshift=5.2cm]
    \draw[->,gray!50] (-1.9,0) -- (1.9,0) node[right,font=\tiny]{$\mathrm{Re}$};
    \draw[->,gray!50] (0,-1.9) -- (0,1.9) node[above,font=\tiny]{$\mathrm{Im}$};
    \draw[gray!30,dashed] (0,0) circle (1);
    \draw[structurepurple,thick,domain=0:360,samples=160,smooth,variable=\deg]
      plot ({(1 + 0.8*cos(\deg))*cos(\deg)}, {-(1 + 0.8*cos(\deg))*sin(\deg)});
    \draw[->,specred,thick] (0,0) -- (0.4,0);
    \fill[specred] (0.4,0) circle (2.2pt) node[above right=1pt,font=\footnotesize\bfseries]{$\mu = (0.4, 0)$};
    \node[anchor=north,font=\small] at (0,-2.1) {(b) $f_{\text{test}} = 3$\,Hz $= f_{\text{signal}}$};
    \node[anchor=north,font=\scriptsize,gray!70!black] at (0,-2.5) {crests align: mass shifts};
  \end{scope}

  % Panel 3: Center of mass distance vs winding frequency
  \begin{scope}[xshift=10.0cm]
    \begin{axis}[
      width=4.8cm, height=4.4cm,
      title={Center of mass vs $f_{\text{test}}$},
      title style={font=\small},
      xlabel={$f_{\text{test}}$ (Hz)}, ylabel={$\|\mu\|$ (strength)},
      xmin=0, xmax=6, ymin=0, ymax=0.6,
      xtick={0,2,3,5}, ytick={0,0.2,0.4},
      axis lines=left,
      tick label style={font=\tiny}, label style={font=\tiny},
    ]
      \addplot[specred, very thick, domain=0.2:5.8, samples=120]
        {abs(0.4*sin(deg(pi*(x-3)))/(pi*(x-3) + 0.001))};
      \addplot[specred, only marks, mark=*, mark size=1.5pt] coordinates {(2,0.02) (3,0.4)};
      \node[font=\tiny, sigblue] at (axis cs:2,0.12) {(a)};
      \node[font=\tiny, structurepurple] at (axis cs:3.4,0.46) {(b)};
    \end{axis}
    \node[anchor=north,font=\small] at (1.8,-2.1) {(c) The frequency spectrum};
    \node[anchor=north,font=\scriptsize,gray!70!black] at (1.8,-2.5) {spike at matching frequency};
  \end{scope}
\end{tikzpicture}
\caption{The Fourier winding machine. To test how much of frequency $f_{\text{test}}$ is present in a signal $x(t)$, wind the waveform around the origin at speed $f_{\text{test}}$ and track its center of mass $\mu = \frac{1}{T}\int_0^T x(t) e^{-\jj 2\pi f_{\text{test}} t}\,\dd t$. (a) When $f_{\text{test}} \neq 3$\,Hz, peaks and troughs disperse symmetrically around the circle and the center of mass stays at the origin. (b) When $f_{\text{test}} = 3$\,Hz, every peak of the wave aligns on the same side, pulling the center of mass outward. (c) Plotting the center-of-mass distance across all winding frequencies reveals the spectrum: an isolated peak exactly at the frequency of the wave.}
\label{ch05:fig-winding}
\end{figure}
\section{Harmonics and timbre: why a violin is not a flute}
When a frequency-$f_0$ tone and its integer multiples $2f_0, 3f_0, \dots$ all
appear together, $f_0$ is the \emph{fundamental} and the multiples are
\emph{harmonics} (or overtones). The ear assigns the \emph{pitch} from the
fundamental, but the \emph{relative amplitudes of the harmonics} are what make a
trumpet sound like a trumpet and a clarinet like a clarinet, even on the same note.
That recipe of harmonic strengths is the \emph{timbre}, sketched for two cases in
Figure~\ref{ch05:timbre}.
\begin{figure}[H]
\centering
\begin{tikzpicture}
\begin{axis}[
  width=12cm, height=4.6cm,
  xlabel={harmonic number $k$}, ylabel={relative amplitude},
  ybar, bar width=7pt, axis lines=left,
  xmin=0.3, xmax=8.7, ymin=0, ymax=1.1,
  xtick={1,2,3,4,5,6,7,8},
  legend pos=outer north east, legend cell align=left,
  enlarge x limits=0.08,
]
\addplot[fill=sigblue!60, draw=sigblue] coordinates {(1,1) (2,0.5) (3,0.33) (4,0.25) (5,0.2) (6,0.16) (7,0.14) (8,0.12)};
\addlegendentry{``brassy'' (rich highs)}
\addplot[fill=specred!50, draw=specred] coordinates {(1,1) (3,0.33) (5,0.2) (7,0.14)};
\addlegendentry{``hollow'' (odd only)}
\end{axis}
\end{tikzpicture}
\caption{Two timbres on the \emph{same} fundamental. Left bars: a full harmonic
series, energy in every multiple -- bright and brassy. The odd-only series (like an
idealized clarinet's lower register) sounds hollow and woody. Same pitch, different
spectrum, different instrument.}
\label{ch05:timbre}
\end{figure}

\begin{applicationbox}
This is why a synthesizer can imitate instruments by \emph{additive synthesis}:
choose the fundamental for the pitch, then dial in a harmonic recipe for the timbre.
It is also why audio compression works -- the ear barely notices a weak harmonic
sitting next to a strong one (masking), so the encoder can throw the weak one away.
Both tricks are pure spectrum thinking.
\end{applicationbox}

\begin{goalbox}
For the wrist controller, ``timbre'' has a motion analogue. A deliberate finger tap
and an accidental table bump can have similar peak acceleration in time, but their
\emph{harmonic content} differs: a crisp tap is a sharp transient (broadband, lots
of highs), while a soft bump is smoother (energy concentrated low). Analyzing the
shape of the spectrum alongside the peak height provides a direct, computationally efficient way to
tell a real gesture from noise. We will turn this into real features in
Chapter~38.
\end{goalbox}

\section{Beats: two close frequencies}

Add two sinusoids of nearly equal frequency and something striking happens: a slow
throbbing in loudness. This is \emph{beating}, and the algebra is a one-line
revelation. With $f_1 = f_0 - \tfrac{\Delta}{2}$ and $f_2 = f_0 + \tfrac{\Delta}{2}$,
\[
  \cos(2\pi f_1 t) + \cos(2\pi f_2 t)
  = \underbrace{2\cos(\pi \Delta\, t)}_{\text{slow envelope}}\;
    \underbrace{\cos(2\pi f_0 t)}_{\text{fast carrier}} .
\]
The sum equals a fast tone at the \emph{average} frequency $f_0$, multiplied by a
slow envelope that wobbles at the \emph{difference} frequency $\Delta$. You hear the
average pitch swelling and fading $\Delta$ times a second. Figure~\ref{ch05:beats}
plots the result.

\begin{figure}[H]
\centering
\begin{tikzpicture}
\begin{axis}[
  width=12cm, height=5cm,
  title={Beats: envelope modulation at difference frequency $\Delta = 3$\,Hz},
  title style={font=\small},
  xlabel={time $t$ (s)}, ylabel={amplitude},
  domain=0:1, samples=200, axis lines=left,
  xmin=0, xmax=1, ymin=-2.3, ymax=2.3,
]
\addplot[sigblue, thick] {cos(deg(2*pi*20*x)) + cos(deg(2*pi*23*x))};
% envelope
\addplot[specred, very thick, dashed] {2*cos(deg(pi*3*x))};
\addplot[specred, very thick, dashed, forget plot] {-2*cos(deg(pi*3*x))};
\end{axis}
\end{tikzpicture}
\caption{Beats: $\cos(2\pi\cdot 20 t)+\cos(2\pi\cdot 23 t)$ in blue. The dashed red
envelope wobbles at the difference $\Delta = 3$~Hz, so the loudness swells three
times a second while the underlying pitch sits at the average $21.5$~Hz.}
\label{ch05:beats}
\end{figure}

\begin{historybox}
The intellectual journey of frequency stretches from ancient harmony to Victorian acoustic
detectors. Around 530~BCE, \textbf{Pythagoras of Samos} conducted the first quantitative
experiment in physics using the monochord: a single string stretched over a movable bridge.
He discovered that stopping the string at exact small-integer length ratios ($2:1$ for the octave,
$3:2$ for the fifth, $4:3$ for the fourth) produced sweet consonant intervals, while complex ratios
produced jarring discord. For two millennia, this remained a mystical numerical dogma of
``universal harmony.''

The physical breakthrough came in 1863 with \textbf{Hermann von Helmholtz} in his monumental treatise
\emph{On the Sensations of Tone}. Long before electronic microphones, amplifiers, or oscilloscopes
existed, Helmholtz asked a profound question: how can we prove that a complex musical sound -- like
a sung vowel, a violin note, or a church bell -- is physically composed of a stack of sinusoidal
harmonics? Helmholtz built a set of tuned spherical glass and brass cavities, now called
\emph{Helmholtz resonators}. Each sphere had two necks: a wide funnel to capture room acoustic
pressure, and a narrow nipple inserted directly into the researcher's ear canal. Each cavity
resonated mechanically at a single, precisely tuned acoustic frequency, heavily amplifying that
single harmonic while leaving all others silent. By placing resonators of sequentially higher
pitches into his ear while an instrument sounded, Helmholtz physically isolated and heard each
individual harmonic ringing by itself! He demonstrated experimentally that what the human brain
perceives as \emph{timbre} is nothing other than the Fourier spectrum: the recipe of harmonic
magnitudes.

Helmholtz's work triggered a famous nineteenth-century controversy with acoustic instrument builder
\textbf{Rudolph K\"{o}nig} over \emph{Ohm's acoustical law} (proposed by Georg Ohm in 1843).
K\"{o}nig insisted that the ear could detect changes in phase, while Helmholtz maintained that the
ear acts as a phase-deaf power-spectrum analyzer. Modern psychoacoustics reveals that while the
cochlea's basilar membrane performs a mechanical Fourier breakdown that is largely phase-insensitive
for sustained vowels, phase alignment critically governs waveform crest factors, speech intelligibility,
and transient perception.
\end{historybox}

\begin{trapbox}
Beating is \emph{not} a new frequency. The spectrum of the beating signal has exactly
two lines, at $20$ and $23$~Hz -- nothing at $3$~Hz and nothing at $21.5$~Hz. The
throb at $3$~Hz is what your \emph{loudness perception} (a nonlinear thing) makes of
two pure tones. The lesson: an envelope you see in time need not be a line in the
spectrum. Time-domain appearance and frequency-domain content are different views,
and reading one off the other takes care.
\end{trapbox}

\section{Complex exponentials and the meaning of negative frequency}

Real cosines are clumsy to compute with because the angle-sum identities multiply
like rabbits. The fix, met in Chapter~4, is Euler's formula:
\[
  e^{\jj\theta} = \cos\theta + \jj\sin\theta .
\]
A complex exponential $e^{\jj 2\pi f t}$ is a point on the unit circle in the complex
plane, spinning counterclockwise at $f$ revolutions per second. Its real part traces
a cosine, its imaginary part a sine. One spinning arrow -- a \emph{phasor} -- carries
the cosine and sine together. This is not a notational convenience; it is the cleaner
truth, because the rotation \emph{is} the frequency, made visible as motion.
Figure~\ref{ch05:phasors} shows the two senses of rotation side by side.

\begin{figure}[H]
\centering
\begin{tikzpicture}[scale=1.5]
  \draw[->, gray] (-1.3,0) -- (1.3,0) node[right] {Re};
  \draw[->, gray] (0,-1.3) -- (0,1.3) node[above] {Im};
  \draw[gray!50] (0,0) circle (1);
  % positive freq phasor
  \draw[->, sigblue, very thick] (0,0) -- (0.707,0.707)
        node[midway, above left, sigblue] {$e^{+\jj\omega t}$};
  \draw[sigblue, ->] (0.95,0.31) arc (18:72:1);
  \node[sigblue] at (1.15,0.75) {CCW};
  % negative freq phasor
  \draw[->, specred, very thick] (0,0) -- (0.707,-0.707)
        node[midway, below left, specred] {$e^{-\jj\omega t}$};
  \draw[specred, ->] (0.95,-0.31) arc (-18:-72:1);
  \node[specred] at (1.15,-0.78) {CW};
\end{tikzpicture}
\caption{Two phasors at $t>0$. The blue arrow $e^{+\jj\omega t}$ spins
counterclockwise -- positive frequency. The red arrow $e^{-\jj\omega t}$ spins
clockwise -- negative frequency. Same speed, opposite sense of rotation.}
\label{ch05:phasors}
\end{figure}

So what on earth is a \emph{negative} frequency? Geometrically it is the cleanest
thing in the world: \textbf{negative frequency just means the phasor spins the other
way.} Counterclockwise is positive; clockwise is negative. There is nothing
mysterious or unphysical about turning the other direction.

\begin{whybox}
Why negative frequencies are forced on us. Rearrange Euler's formula and you get the
identity at the heart of all real-signal spectra:
\[
  \cos(2\pi f t) = \tfrac12 e^{\jj 2\pi f t} + \tfrac12 e^{-\jj 2\pi f t}.
\]
A real cosine is the \emph{average of two counter-rotating phasors}, one at $+f$ and
one at $-f$. Picture them: two arrows spinning in opposite directions. Their vertical
(imaginary) parts always cancel, so the sum is purely real; their horizontal parts
reinforce, tracing the cosine. To build a real wave from complex atoms you
\emph{must} use both senses of rotation. That is why the spectrum of any real signal
is symmetric: every component at $+f$ has a mirror twin at $-f$. Negative frequency
is not a fiction; it is the bookkeeping that lets one elegant complex framework
describe real signals.
\end{whybox}

\begin{happeningbox}
The cancellation is worth seeing slowly. At each instant the $+f$ phasor is at angle
$+\omega t$ and the $-f$ phasor at angle $-\omega t$. They are mirror images across
the real axis. Add two mirror-image arrows: the up-down parts kill each other, the
left-right parts double. The result slides back and forth along the real axis -- a
cosine. Frequency $+f$ and frequency $-f$ are the two ``hands'' that together clap
out a real oscillation.
\end{happeningbox}

\begin{focusbox}
For \emph{real} signals, negative frequencies are redundant: knowing the $+f$ half
tells you the $-f$ half by symmetry, so many textbooks (and tuners and audio plots)
quietly fold them away and show only $f \ge 0$. But for \emph{complex} signals -- the
I/Q data of a radio, or the analytic signal we build in Chapter~25 -- positive and
negative frequencies are genuinely different and independent. A single phasor
spinning one way has energy at exactly one frequency, positive or negative, with no
twin. That asymmetry is the whole basis of how a radio tells an upper sideband from a
lower one.
\end{focusbox}

\begin{connectionbox}
The probability companion will lean on this same complex-exponential atom: the
\emph{characteristic function} of a random variable is $\E[e^{\jj\omega X}]$, the
expected value of a phasor whose spin rate is set by the data. It is the Fourier
transform of a probability density, and convolving two densities (adding two
independent random variables) becomes multiplying their characteristic functions --
the very same ``convolution becomes multiplication'' magic that makes frequency
thinking pay off here. One idea, two books.
\end{connectionbox}

\section{From a few lines to a continuous spectrum}

So far every example was \emph{periodic}, and its spectrum was a discrete set of
lines at the fundamental and its harmonics. What about a signal that does not repeat
-- a single finger tap, a click, a one-shot transient?

The bridge is to imagine stretching the period $T$ toward infinity. The harmonics
sit at multiples of $f_0 = 1/T$; as $T$ grows, those allowed frequencies crowd closer
and closer together (spacing $1/T \to 0$). In the limit the discrete comb of lines
fuses into a continuous curve. A non-periodic signal has a \emph{continuous} spectrum:
energy spread over a range of frequencies rather than parked on a few exact ones.
That limit is the Fourier \emph{transform}, the subject of Chapter~7; the line-spectrum
case is the Fourier \emph{series} of Chapter~6. Figure~\ref{ch05:disc-to-cont}
contrasts the two.

\begin{figure}[H]
\centering
\begin{tikzpicture}
% discrete
\begin{axis}[
  name=disc, width=6cm, height=4cm,
  title={Periodic: discrete lines},
  title style={font=\small},
  xlabel={$f$}, ylabel={amplitude},
  axis lines=left, xmin=0, xmax=6, ymin=0, ymax=1.2,
  xtick={1,2,3,4,5}, ycomb, mark=*, mark options={sigblue},
]
\addplot[sigblue, very thick] coordinates {(1,1) (2,0.5) (3,0.33) (4,0.25) (5,0.2)};
\end{axis}
% continuous
\begin{axis}[
  at={(disc.east)}, anchor=west, xshift=1.3cm,
  width=6cm, height=4cm,
  title={One-shot: continuous curve},
  title style={font=\small},
  xlabel={$f$}, ylabel={magnitude},
  axis lines=left, domain=0:6, samples=120,
  xmin=0, xmax=6, ymin=0, ymax=1.2,
]
\addplot[specred, very thick] {abs(sin(deg(pi*x*0.5))/(pi*x*0.5 + 0.0001))};
\end{axis}
\end{tikzpicture}
\caption{Left: a periodic signal has a discrete line spectrum (energy only at
harmonics). Right: a non-repeating pulse has a continuous spectrum (a smooth curve --
here a sinc-like magnitude). Stretching the period to infinity merges the lines into
the curve.}
\label{ch05:disc-to-cont}
\end{figure}

\begin{intuitionbox}
A short, sharp event must be \emph{broadband}: to localize energy tightly in time you
need many frequencies cooperating, just as it took many harmonics to make the square
wave's sharp edge. A long, pure tone is \emph{narrowband}: one frequency, spread over
a long time. Brief in time $\Leftrightarrow$ wide in frequency. This trade-off is the
seed of the time--frequency uncertainty principle (Chapter~10) and the reason a drum
(brief) sounds ``noisy'' while an organ pipe (sustained) sounds ``pure''.
\end{intuitionbox}

\section{A worked example end to end}

\begin{example}[Reading a signal's spectrum by inspection]
\label{ch05:ex-inspect}
Consider
\[
  x(t) = 3 + 2\cos(2\pi\cdot 10\,t) + \cos(2\pi\cdot 60\,t - \tfrac{\pi}{2}).
\]
Describe its spectrum without any integral.

\textbf{The constant $3$.} A constant is a cosine of frequency zero ($\cos(0)=1$),
amplitude $3$. This is the \emph{DC component} -- the signal's average value. It is a
line at $f=0$ of height $3$.

\textbf{The $2\cos(2\pi\cdot 10 t)$ term.} Amplitude $2$, frequency $10$~Hz, phase
$0$. A line at $10$~Hz of height $2$ (and, by symmetry, a twin at $-10$~Hz if we use
the two-sided complex picture).

\textbf{The $\cos(2\pi\cdot 60 t - \pi/2)$ term.} Amplitude $1$, frequency $60$~Hz,
phase $-\pi/2$ (which makes it a \emph{sine}, $\sin(2\pi\cdot 60 t)$). A line at
$60$~Hz of height $1$, carrying a phase of $-\pi/2$.

\textbf{Spectrum (one-sided magnitude):} spikes at $0, 10, 60$~Hz of heights
$3, 2, 1$. Nothing anywhere else. We read off the complete frequency content of a
signal that, drawn in time, looks like a complicated wiggle -- and we never
integrated. That is the fluency this chapter is for. Figure~\ref{ch05:ex-spectrum}
draws the three lines.
\end{example}

\begin{figure}[H]
\centering
\begin{tikzpicture}
\begin{axis}[
  width=11cm, height=4.4cm,
  title={Magnitude spectrum of $x(t)$ (three spectral lines)},
  title style={font=\small},
  xlabel={frequency $f$ (Hz)}, ylabel={magnitude},
  axis lines=left, xmin=-5, xmax=70, ymin=0, ymax=3.5,
  xtick={0,10,60}, ytick={0,1,2,3},
  ycomb, mark=*, mark options={specred},
]
\addplot[specred, very thick] coordinates {(0,3) (10,2) (60,1)};
\end{axis}
\end{tikzpicture}
\caption{One-sided magnitude spectrum of
$x(t)=3+2\cos(2\pi\cdot 10\,t)+\cos(2\pi\cdot 60\,t-\pi/2)$. A DC line of height
$3$, the $10$~Hz tone of height $2$, and the $60$~Hz tone of height $1$ -- read off
by inspection.}
\label{ch05:ex-spectrum}
\end{figure}

\begin{walkbox}
Suppose the $60$~Hz term were unwanted hum from the wall socket contaminating a
sensor. Knowing it lives \emph{only} at $60$~Hz, how would you remove it without
harming the $10$~Hz signal you care about? (Answer foreshadowed: a filter that passes
$10$~Hz and rejects $60$~Hz -- a notch or low-pass. The entire reason filtering is
possible is that the interference and the signal occupy different frequencies. We
build such filters in Chapter~16.)
\end{walkbox}

\section{Frequency for the wrist controller}

\begin{goalbox}
Let us anchor every idea above in the device. Typical frequency bands you will meet
once the IMU and EMG are strapped on:
\begin{itemize}
  \item \textbf{Gravity / posture}: essentially DC up to $\sim 0.5$~Hz on the
        accelerometer. Slow, low frequency -- the gross orientation of the forearm.
  \item \textbf{Deliberate gestures}: roughly $1$--$8$~Hz. A comfortable finger-tap
        or wrist-flick rhythm lives here. This is the band we actually want.
  \item \textbf{Tremor}: a physiological hand tremor sits around $8$--$12$~Hz -- close
        enough to fast gestures that telling them apart is a real spectral problem.
  \item \textbf{Muscle activity (EMG)}: surface EMG energy spreads broadly, roughly
        $20$--$450$~Hz, which is exactly why EMG needs a much higher sampling rate
        than the IMU.
  \item \textbf{Mains hum}: a sharp interfering line at $50$ or $60$~Hz (and its
        harmonics), right in the middle of the EMG band -- a notch filter's classic
        target.
\end{itemize}
Every design choice -- sampling rate, filter cutoffs, which features to compute --
flows from where in frequency the things you care about and the things you do not
care about happen to sit. Figure~\ref{ch05:bands} lays the bands out on one axis.
\end{goalbox}

\begin{figure}[H]
\centering
\begin{tikzpicture}
\begin{semilogxaxis}[
  width=13cm, height=4.6cm,
  title={Frequency bands for wearable sensing and gesture control},
  title style={font=\small},
  xlabel={frequency (Hz, log scale)},
  axis lines=left, xmin=0.1, xmax=600,
  ymin=0, ymax=4.2, ytick=\empty,
  xtick={0.1,1,10,100},
  xticklabels={$0.1$,$1$,$10$,$100$},
]
% gravity/posture: 0.1 to 0.5
\addplot[draw=none, fill=sigblue!35] coordinates {(0.1,0.4)(0.5,0.4)(0.5,0.9)(0.1,0.9)} \closedcycle;
\node[anchor=west, sigblue, font=\footnotesize] at (axis cs:0.105,1.3) {gravity / posture};
% gestures: 1 to 8
\addplot[draw=none, fill=targetgreen!45] coordinates {(1,1.6)(8,1.6)(8,2.1)(1,2.1)} \closedcycle;
\node[anchor=west, targetgreen!60!black, font=\footnotesize] at (axis cs:1,2.45) {gestures};
% tremor: 8 to 12
\addplot[draw=none, fill=quotientorange!55] coordinates {(8,1.6)(12,1.6)(12,2.1)(8,2.1)} \closedcycle;
\node[anchor=west, quotientorange, font=\footnotesize] at (axis cs:9,2.8) {tremor};
% EMG: 20 to 450
\addplot[draw=none, fill=structurepurple!35] coordinates {(20,0.4)(450,0.4)(450,0.9)(20,0.9)} \closedcycle;
\node[anchor=west, structurepurple, font=\footnotesize] at (axis cs:22,1.3) {muscle (EMG)};
% mains hum line at 60
\addplot[specred, very thick] coordinates {(60,0)(60,3.6)};
\node[anchor=south, specred, font=\footnotesize] at (axis cs:60,3.6) {$60$~Hz hum};
\end{semilogxaxis}
\end{tikzpicture}
\caption{Where the wrist controller's signals live in frequency. Slow posture sits
near DC; deliberate gestures occupy $1$--$8$~Hz with tremor pressing right against
them at $8$--$12$~Hz; surface EMG spreads broadly to a few hundred hertz, with the
mains hum a sharp line at $60$~Hz cutting through it.}
\label{ch05:bands}
\end{figure}

\begin{buildbox}
A first hands-on frequency check, before any fancy transform. Sample the
accelerometer steadily, then estimate the rate of a repetitive gesture by counting
zero-crossings of the (mean-removed) signal: a sinusoid at $f$~Hz crosses zero $2f$
times per second.
\begin{lstlisting}[language=Python]
import numpy as np
def gesture_rate_hz(x, fs):
    x = x - np.mean(x)                      # remove DC (gravity/offset)
    crossings = np.sum((x[:-1] < 0) & (x[1:] >= 0))  # rising zero-crossings
    duration = len(x) / fs
    return crossings / duration             # one crossing per cycle
# fs = 100 Hz IMU stream; x = one window of one axis
\end{lstlisting}
Crude, fooled by noise, and blind to anything but the dominant rhythm -- but it runs
in microcontroller-friendly integer-ish math and gives a usable tap-tempo estimate
before you have an FFT. We graduate to the real spectrum in Chapter~8.
\end{buildbox}

\begin{trapbox}
\textbf{What a wrist IMU can and cannot do.}
It is tempting to imagine the spectrum of a wrist accelerometer revealing individual
finger movements. It cannot. A wrist-mounted IMU measures the gross translation and
rotation of the \emph{wrist}; the tiny, mechanically isolated motions of individual
finger joints barely couple into it. The IMU can tell ``a tap happened'' or ``the
wrist rolled'', i.e.\ coarse, discrete gestures -- it cannot recover continuous,
per-finger joint angles. Distinguishing which finger tapped, or tracking finger pose,
needs \emph{other physics}: surface EMG or force myography reading the forearm
muscles that drive the fingers (Chapter~37), or sensors actually on the fingers
(instrumented rings/gloves). Frequency analysis sharpens whatever the sensor can
sense; it cannot conjure information the sensor never captured. We will keep saying
this, because overclaiming here is the most common way these projects fail.
\end{trapbox}

\section{What to carry forward}

\begin{recoverybox}
\begin{itemize}
  \item Frequency is a \emph{rate of repetition}. Three names for it: period $T$,
        ordinary frequency $f=1/T$ in Hz, angular frequency $\omega=2\pi f$ in rad/s.
  \item A pure sinusoid is fixed by three numbers: amplitude, frequency, phase.
  \item Signals can be \emph{built} from sinusoids. Low frequencies carry gross
        shape; high frequencies carry sharp detail and edges.
  \item The \emph{spectrum} lists how much of each frequency is present (magnitude)
        and its timing (phase). For periodic signals it is a few lines; for
        one-shot signals it is a continuous curve.
  \item Harmonics and their relative strengths are \emph{timbre}; beats are two close
        frequencies producing a slow loudness throb (and are \emph{not} a new
        spectral line).
  \item A complex exponential is a spinning phasor; \emph{negative frequency} just
        means spinning the other way. Real signals need both $+f$ and $-f$ twins,
        which is why real spectra are symmetric.
  \item Brief in time $\Leftrightarrow$ broad in frequency.
\end{itemize}
Next chapter we make the periodic case exact: the Fourier series, where the harmonic
amplitudes stop being guesses and become inner products you can compute.
\end{recoverybox}

\begin{connectionbox}
\textbf{Capabilities, Downsides, and Forward Roadmap}
\begin{itemize}
  \item \textbf{What this unlocked}: The geometric and physical concept of frequency as repetition,
  winding, and resonance. We can now explain timbre through harmonic distributions, visualize
  negative frequency as counter-rotation, understand how the ear mechanically splits sound like a bank
  of Helmholtz resonators, and intuitively view Fourier analysis as finding the center of mass of a
  wound wire in the complex plane.
  \item \textbf{Honest downsides and limitations}: The intuitive winding-machine model is qualitative;
  it does not tell us the exact weighting factors needed to reconstruct the original continuous signal,
  nor how to prevent spectral leakage when winding non-integer periods. Furthermore, continuously
  winding a signal around the origin for every arbitrary frequency $\omega$ is computationally
  intractable without an organized basis of orthogonal functions.
  \item \textbf{Where we go next}: Chapter 6 takes the winding machine and formalizes it into the
  exact mathematics of the Fourier Series for strictly periodic signals. Chapter 7 stretches the
  period to infinity ($T_0 \to \infty$) to create the Continuous Fourier Transform for transient
  signals; and Chapter 8 adapts the winding machine into finite matrix operations for discrete
  signals on digital microprocessors.
\end{itemize}
\end{connectionbox}

\section{Exercises}

\begin{exercise}
A signal repeats every $40$~ms. Give its period in seconds, its ordinary frequency in
hertz, and its angular frequency in radians per second.
\end{exercise}

\begin{exercise}
Write $5\cos(120\pi t + \pi/3)$ in the form $A\cos(2\pi f t + \phi)$, identifying
$A$, $f$ (Hz), and $\phi$. Is $120\pi$ an angular frequency or an ordinary one? How
do you know?
\end{exercise}

\begin{exercise}
Using $\cos\theta = \tfrac12(e^{\jj\theta}+e^{-\jj\theta})$, write
$x(t)=4\cos(2\pi\cdot 50\,t)$ as a sum of two complex exponentials. At which two
frequencies (including sign) does it have energy, and what is the magnitude at each?
\end{exercise}

\begin{exercise}
Two tuning forks sound together and you count $5$ loudness swells per second. One fork
is known to be exactly $440$~Hz. What are the two possible frequencies of the other
fork? Explain why two answers are possible from the beat rate alone.
\end{exercise}

\begin{exercise}
Sketch (by reasoning, not computing) the magnitude spectrum of
$x(t) = 1 + \cos(2\pi\cdot 5 t) + \tfrac12\sin(2\pi\cdot 5 t)$. Hint: combine the two
$5$~Hz terms into a single cosine with a phase, then list the lines.
\end{exercise}

\begin{exercise}
A square wave of fundamental $f_0 = 100$~Hz is built from odd harmonics with
amplitudes $\propto 1/k$. List the frequencies and relative amplitudes of its first
four nonzero spectral lines.
\end{exercise}

\begin{exercise}
Explain in one or two sentences each: (a) why a perfectly sharp step edge in a signal
requires arbitrarily high frequencies; (b) why a pure, long-sustained tone occupies
essentially a single frequency.
\end{exercise}

\begin{exercise}
A real signal has spectral lines at $\pm 30$~Hz and $\pm 90$~Hz. Argue from symmetry
that the signal is real, and write a plausible time-domain expression that produces
exactly these lines.
\end{exercise}

\begin{exercise}
The zero-crossing rate estimator in the build box reports $8$ rising crossings in a
$2$-second accelerometer window after mean removal. What gesture frequency does it
estimate? Name two situations (think noise and harmonics) where this estimate would
be wrong, and say in which direction.
\end{exercise}

\begin{exercise}
You record forearm EMG and see a strong, narrow $60$~Hz line plus broadband energy
from $20$ to $300$~Hz. Which part is the muscle signal and which is interference?
Roughly where in frequency would you place a filter to suppress the interference
while keeping most of the muscle energy, and what kind of filter is it?
\end{exercise}

\begin{exercise}
(Conceptual.) A friend claims that because they can clearly see a slow throbbing
envelope at $3$~Hz when two $1$~kHz-ish tones beat together, ``there must be a
$3$~Hz component in the signal.'' Use the spectrum of the beating signal to explain
exactly what is and is not present at $3$~Hz.
\end{exercise}

\begin{exercise}
(Toward the build.) Argue, using the limits discussed in this chapter, why no
amount of clever frequency analysis of a single \emph{wrist} accelerometer can
recover which individual finger tapped. What additional sensing would you add, and
what physical quantity would it measure?
\end{exercise}
